\documentclass[12pt,a4paper]{article}

\usepackage[T1]{fontenc}
\usepackage[utf8]{inputenc}
\usepackage[brazil]{babel}
\usepackage{lmodern}
\usepackage{geometry}
\usepackage{hyperref}
\usepackage{enumitem}
\usepackage{array}

\geometry{margin=2.5cm}
\hypersetup{
  colorlinks=true,
  linkcolor=blue,
  urlcolor=blue
}

\title{Plano de Aula: Aula 6 -- Modelagem de Software (UML)}
\date{}

\begin{document}
\maketitle

\noindent
\textbf{Disciplina:} Engenharia de Software (Curso Técnico em Desenvolvimento de Sistemas)\\
\textbf{Duração:} 2 horas (120 minutos)\\
\textbf{Modalidade:} Presencial\\
\textbf{Materiais da aula:} \href{aula-06.html}{slides (HTML)} \quad \href{aula-06.pdf}{ementa (PDF)}\\
\textbf{Livro-base (referência):} \href{../99-Materiais/Engenharia%20de%20software%20projetos%20e%20processos%20by%20Wilson%20de%20P%C3%A1dua%20Paula%20Filho%20.epub.pdf}{Engenharia de software: projetos e processos (PDF)}

\section{Competências e Habilidades}
\textbf{Competências}
\begin{itemize}[leftmargin=*,itemsep=0.25em]
  \item Compreender e aplicar modelagem para apoiar análise e projeto.
  \item Analisar e projetar softwares orientados a objetos.
\end{itemize}

\textbf{Habilidades}
\begin{itemize}[leftmargin=*,itemsep=0.25em]
  \item Conhecer e utilizar UML como linguagem de especificação.
  \item Representar requisitos (casos de uso), estrutura (classes) e comportamento (sequência).
  \item Manter consistência entre modelos e implementação.
\end{itemize}

\section{Objetivos da Aula}
\begin{itemize}[leftmargin=*,itemsep=0.25em]
  \item Explicar por que modelar e quando modelar ``o suficiente''.
  \item Produzir um caso de uso simples com ator, objetivo e fluxo.
  \item Produzir diagrama de classes básico do domínio.
  \item Produzir um diagrama de sequência para o fluxo principal.
\end{itemize}

\section{Assuntos Abordados no Dia}
\begin{itemize}[leftmargin=*,itemsep=0.25em]
  \item Modelagem como comunicação e redução de risco.
  \item UML e perspectivas: requisitos, estrutura e comportamento.
  \item Caso de uso: atores, objetivos, fluxo principal e alternativas.
  \item Classes: conceitos do domínio, atributos, métodos e relações.
  \item Sequência: mensagens e responsabilidades no tempo.
\end{itemize}

\section{Metodologia}
Aula expositiva com demonstrações no quadro e exercícios guiados. Os estudantes desenvolvem modelos em grupos a partir do ``app organizador'' e apresentam rapidamente para checagem de consistência. O professor enfatiza rastreabilidade: requisito $\rightarrow$ modelo $\rightarrow$ implementação $\rightarrow$ testes.

\section{Materiais e Recursos}
\begin{itemize}[leftmargin=*,itemsep=0.25em]
  \item Quadro e marcadores.
  \item Papel A4 e canetas para desenhar diagramas (ou ferramenta CASE, se disponível).
  \item Slides: \href{aula-06.html}{aula-06.html} e ementa: \href{aula-06.pdf}{aula-06.pdf}.
  \item Livro-base (PDF): \href{../99-Materiais/Engenharia%20de%20software%20projetos%20e%20processos%20by%20Wilson%20de%20P%C3%A1dua%20Paula%20Filho%20.epub.pdf}{Engenharia de software: projetos e processos}.
  \item Vídeos complementares (links nos slides).
\end{itemize}

\section{Cronograma e Descrição Detalhada (120 Minutos)}
\subsection*{Parte 1: Abertura (10 min)}
\begin{itemize}[leftmargin=*,itemsep=0.35em]
  \item Conectar com prototipagem (Aula 5): protótipo valida interface; modelo valida estrutura e regras.
  \item Apresentar objetivos e entregas: caso de uso, classes e sequência.
\end{itemize}

\subsection*{Parte 2: UML e ``quanto modelar'' (20 min)}
\begin{itemize}[leftmargin=*,itemsep=0.35em]
  \item Mostrar perspectivas e quando usar cada diagrama.
  \item Discutir risco de excesso: documentação que não vira produto.
  \item Regra prática: modelar para decidir e construir com segurança.
\end{itemize}

\subsection*{Parte 3: Caso de uso (20 min)}
\begin{itemize}[leftmargin=*,itemsep=0.35em]
  \item Explicar estrutura: ator, objetivo, fluxo principal, alternativas.
  \item Exercício guiado: ``Cadastrar tarefa'' (fluxo principal + 1 alternativa).
\end{itemize}

\subsection*{Parte 4: Classes (25 min)}
\begin{itemize}[leftmargin=*,itemsep=0.35em]
  \item Identificar classes do domínio do ``app organizador''.
  \item Definir atributos mínimos e relacionamentos.
  \item Discutir responsabilidades e evitar ``classe Deus''.
\end{itemize}

\subsection*{Parte 5: Sequência (25 min)}
\begin{itemize}[leftmargin=*,itemsep=0.35em]
  \item Montar sequência do fluxo principal do caso de uso.
  \item Discutir responsabilidades e dependências.
\end{itemize}

\subsection*{Parte 6: Apresentação e fechamento (20 min)}
\begin{itemize}[leftmargin=*,itemsep=0.35em]
  \item Grupos apresentam modelos (2--3 min).
  \item Checagem de consistência: caso de uso $\leftrightarrow$ classes $\leftrightarrow$ sequência.
  \item Ponte para Aula 7: como modelos ajudam a derivar casos de teste.
\end{itemize}

\section{Avaliação}
\begin{itemize}[leftmargin=*,itemsep=0.25em]
  \item Entrega do conjunto de modelos (clareza e consistência).
  \item Participação na discussão e revisão com feedback.
\end{itemize}

\section{Referências}
\begin{itemize}[leftmargin=*,itemsep=0.25em]
  \item \href{../99-Materiais/Engenharia%20de%20software%20projetos%20e%20processos%20by%20Wilson%20de%20P%C3%A1dua%20Paula%20Filho%20.epub.pdf}{PAULA FILHO, Wilson de Pádua. Engenharia de software: projetos e processos. 4. ed. 2019.}
\end{itemize}

\end{document}
